\documentclass[11pt,a4paper]{article}
\usepackage[T1]{fontenc}
\usepackage[margin=25mm]{geometry}
\usepackage{amsmath,amssymb,amsthm,booktabs}
\usepackage[hidelinks]{hyperref}
\newtheorem{proposition}{Proposition}
\newtheorem{lemma}{Lemma}
\title{A projective-line code disproves conjecture 00000007973}
\author{asfigh665}
\date{October 15, 2026}
\begin{document}
\maketitle
\begin{abstract}
The conjectured minimum-weight identity $d=\deg D-g$ fails for a standard
evaluation algebraic-geometric code on $\mathbf P^1_{\mathbf F_3}$.
For $D=P_0+P_1+P_2$ and $G=P_\infty$, the code is the $[3,2,2]_3$
Reed--Solomon code, whereas the proposed identity gives $d=3$.
A dependency-free Lean 4 project checks the finite field, the complete
evaluation code, its minimum nonzero weight, and its weight enumerator.
\end{abstract}

\section{The assertion being refuted}
Conjecture 00000007973 concerns the weight enumerator of a code
$C(X,D,G)$ from algebraic geometry. Both language versions explicitly
assert that its minimum weight is $\deg D-g$, where $g$ is the genus of $X$.
As usual, minimum weight means the minimum over \emph{nonzero} codewords.
We disprove this clause; consequently the conjunction of assertions in the
conjecture is false, independently of its proposed modular-form description.

We use the standard evaluation AG code
\[
 C_L(D,G)=\{(f(P_0),f(P_1),f(P_2)):f\in L(G)\},
 \qquad L(G)=\{0\}\cup\{f\in\mathbf F_3(X)^\times:(f)+G\geq0\}.
\]
Here $X=\mathbf P^1_{\mathbf F_3}$, $t$ is its affine coordinate,
$P_j$ is the rational point $t=j$, $D=P_0+P_1+P_2$, and $G=P_\infty$.
The curve is smooth, projective, and geometrically integral, with genus
$g=0$. The evaluation points are distinct, and the supports of $D$ and
$G$ are disjoint. Also $0<\deg G=1<\deg D=3$ and $2g-2<\deg G$,
so the example satisfies the usual nondegenerate evaluation-code hypotheses.

\begin{lemma}
$L(P_\infty)=\{at+b:a,b\in\mathbf F_3\}$.
\end{lemma}
\begin{proof}
A rational function on the projective line with no finite poles is a
polynomial in $t$. Indeed, write it as $p(t)/q(t)$ in lowest terms; any
nonconstant irreducible factor of $q$ supplies a finite closed point where
the function has a pole. Thus $q$ must be constant. In the local parameter
$u=1/t$ at infinity, a nonzero polynomial of degree $m$ has pole order $m$.
The bound of one on that pole order is therefore exactly $\deg f\leq1$.
The zero function is included separately.
\end{proof}
In particular the full Weierstrass semigroup at $P_\infty$ is
$\mathbb N_0$: the function $t^m$ has its unique pole there, of order $m$.
Its gap set is empty. Thus the semigroup data mentioned in the conjecture
are completely specified in this example.

\section{Exact minimum weight}
The lemma gives the whole code explicitly:
\[
 C_L(D,G)=\{(b,a+b,2a+b):a,b\in\mathbf F_3\}.
\]
The parameterization is injective: the value at $0$ recovers $b$, and
the value at $1$ then recovers $a$. Hence this is a length-three,
dimension-two linear code over $\mathbf F_3$.

\begin{proposition}
The minimum nonzero Hamming weight of $C_L(D,G)$ is $2$, and its
weight enumerator is $1+6z^2+2z^3$.
\end{proposition}
\begin{proof}
When $a=b=0$, the word has weight zero. When $a=0$ and $b\ne0$,
the word is a nonzero constant and has weight three; there are two such
words. When $a\ne0$, the polynomial $at+b$ has exactly one root in
$\mathbf F_3$. All three field elements are evaluated, so its word has
exactly two nonzero entries; there are $2\cdot3=6$ such words.
In particular, $f=t$ gives the word $(0,1,2)$ of weight two. These cases
exhaust all nine codewords, proving both claims.
\end{proof}

Thus
\[
 \boxed{d(C_L(D,G))=2\ne3=\deg D-g.}
\]
This discrepancy is not the zero codeword or an exceptional divisor issue.
The familiar designed-distance lower bound is $d\geq\deg D-\deg G$;
it is an equality in this example. The conjecture instead subtracts the
genus, omitting the divisor $G$ on which the evaluation space depends.

\section{Lean verification and its scope}
The accompanying project uses Lean 4.24.0 and its standard library only.
\texttt{F3} is \texttt{Fin 3}, with addition and multiplication modulo three;
\texttt{field\_laws} checks the field identities and multiplicative inverses.
\texttt{evaluate} constructs $t\mapsto at+b$ using those operations.
\texttt{InCode} quantifies over all $a,b$, and
\texttt{MinimumNonzeroWeight} requires both a lower bound for every
nonzero codeword and an attaining codeword. The principal results are:
\begin{itemize}
\item \texttt{evaluation\_linear} and \texttt{evaluate\_injective} verify
the algebraic parameterization;
\item \texttt{minimum\_weight\_two} and
\texttt{not\_minimum\_weight\_three} establish the minimum and refute
the claimed value;
\item \texttt{complete\_weight\_enumerator} checks all nine parameter pairs;
\item \texttt{counterexample} refutes the required genus-zero,
degree-one instance; the final implication theorem explicitly transfers
this contradiction to any assertion implying that instance.
\end{itemize}
All finite decisions are kernel checked. The project does not formalize
schemes or general Riemann--Roch theory; the identification of this concrete
evaluation space with $L(P_\infty)$, and the standard genus of $\mathbf P^1$,
are the mathematical bridge supplied above. No general algebraic-geometry
theorem is asserted as an added axiom in Lean. The checks therefore cover
the full finite code and the contradiction, rather than a hard-coded
comparison of two unexplained integers.

\paragraph{Reproduction.}
From \texttt{lean4/}, run \texttt{lake build},
\texttt{lake env lean Conjecture7973.lean}, and \texttt{lake exe verify}.
The displayed axiom dependencies of the main results are only
\texttt{propext} and \texttt{Quot.sound}. The project has no external packages,
proof holes, native decision shortcuts, or custom axioms.
\end{document}
